%
This search uses two discriminants to identify a possible signal over the large expected background rate from SM processes. 

A first discriminant based on Boosted Decision Trees (BDT)~\ref{sec:mva_bdt} is defined following similar studies performed in the context of the ATLAS measurement of SM production of four-top-quarks~\cite{ATLAS-CONF-2021-013}. The BDT is further optimised for this search as described in the following. A second dedicated discriminant is defined based on a Graph Neural Network~\ref{sec:mva_gnn}. 

The use of the BDT and GNN is combined with a novel structure which incorporates physics parameters of interest, in our case represented by the mass of the heavy Higgs boson \mH, as features for the training~\cite{Baldi_2016}, leading to a so-called parametrised classifier.
Parametrised classifiers provide smooth interpolation across physics parameters and improve the performance of isolated classifiers.
It is trained with signal samples at all mass benchmarks together against \ttbar background.
The training regions are chosen to be $\geq9j\geq3b$ for 1L and $\geq7j\geq3b$ for 2LOS channel.
The yields of signal samples and \ttbar background to train the BDT and GNN are shown in table~\ref{tab:MVAYields}.
The technical implementation is based on introducing a pseudo-mass feature in the background sample used for the training; this is a discrete variable assuming values between 400 to 1000~\gev\ in step of 100~\gev\ as available for the \ttH\ signal benchmarks. 
This pseudo-mass feature is also introduced for signal events as their truth \mH value. The pseudo-mass is treated as an additional input variable in the training. 

Over-training may occur when the performance of both BDT and GNN in the particular set of training samples is deviated from the performance in another set of samples.
It is checked by 2-fold cross-training in this analysis.
The samples are splitted into even and odd, based on the event index.
The output score of the model is evaluated by using the model trained on even events to evaluate odd events and vice-versa.
The area under ROC curve (AUC) is used as a performance metric and the difference between training and testing AUC is used as a metric of overtraining. 

Detailed comparative studies of the two discriminants were performed. The GNN and BDT discriminant modelling in signal depleted phase-spaces (see Sec.~\ref{sec:blinding}) was found to be satisfactory in both cases. The expected analysis sensitivity was observed to be significantly better with the GNN discriminant, leading to 20\%-10\% improvements in the expected exclusion limits depending on the assumed Higgs boson mass and lepton channel in consideration.  

The BDT and GNN discriminants are used both to perform a full analysis. The results for the GNN discriminant are taken as nominal, the results of the BDT analysis are used as cross check.

%%%Table: Yields of Signal and background of MVA Training%%% 
\begin{table}[h!]
\centering
\begin{tabular}{c|c|c|c|c|c|c|c|c}
\hline \hline
\multicolumn{9}{c}{1L Channel} \\
\hline
& $t\bar{t}+$ jets & M400 & M500 & M600 & M700 & M800 & M900 & M1000 \\
Unweighted Yields & 622925 & 26645 & 32370 & 36776 & 40119 & 42436 & 44555 & 46029 \\
Weighted Yields & 4760.97 & 59.90 & 72.34 & 81.59 & 88.60 & 93.29 & 97.8176 & 100.55 \\
\hline
\multicolumn{9}{c}{2LOS Channel} \\
\hline
& $t\bar{t}+$ jets & M400 & M500 & M600 & M700 & M800 & M900 & M1000 \\
Unweighted Yields & 135055 & 8999 & 10561 & 11807 & 12507 & 13225 & 13702 & 14238 \\
Weighted Yields & 1033.87 & 21.09 & 24.58 & 27.18 & 28.78 & 30.17 & 31.12 & 32.27 \\
\hline \hline
\end{tabular}
\caption{\label{tab:MVAYields}The unweighted and weighted yields of signal and \ttbar background used to train BDT and MVA.}
\end{table}

